\chapter{Dasar Teori}

\section{Konsep Dasar Containerization}

\subsection{Docker dan Arsitektur Container}
\textit{Containerization} adalah paradigma virtualisasi tingkat sistem operasi yang memungkinkan aplikasi beserta seluruh dependensinya dikemas dalam satu unit terisolasi yang disebut \textit{container}. Berbeda dengan mesin virtual (\textit{Virtual Machine}) yang menjalankan salinan penuh sistem operasi tamu di atas \textit{hypervisor}, kontainer berbagi \textit{kernel} sistem operasi \textit{host} sehingga menghasilkan \textit{overhead} yang jauh lebih rendah [\ref{felter}]. Studi komparatif oleh Felter et al. menunjukkan bahwa kontainer mampu mencapai kecepatan eksekusi mendekati \textit{native} (\textit{near-native performance}), dengan waktu \textit{startup} dalam hitungan detik dibandingkan mesin virtual yang memerlukan waktu dalam hitungan menit [\ref{felter}].

Docker merupakan platform \textit{containerization} paling dominan yang menyediakan ekosistem lengkap untuk membangun, mendistribusikan, dan menjalankan kontainer. Arsitektur Docker terdiri atas tiga komponen utama: \textit{Docker Engine} (daemon yang mengelola siklus hidup kontainer), \textit{Docker Image} (templat \textit{read-only} berlapis yang mendefinisikan lingkungan aplikasi), dan \textit{Docker Container} (instansi yang dapat dieksekusi dari sebuah \textit{image}) [\ref{dockerdocs}]. Hubungan antara ketiga komponen ini bersifat hierarkis — sebuah \texttt{Dockerfile} menghasilkan \textit{image}, dan \textit{image} diinstansiasi menjadi satu atau lebih kontainer yang berjalan secara terisolasi.

\begin{figure}[H]
    \centering
    \includegraphics[width=0.45\textwidth]{docker-architecture.png}
    \caption{Perbandingan arsitektur containerization versus virtualisasi tradisional}
\end{figure}

\clearpage
\subsection{Dockerfile dan Strategi Layer Caching}
\texttt{Dockerfile} adalah berkas teks yang berisi serangkaian instruksi deklaratif untuk membangun sebuah \textit{Docker image} secara otomatis. Setiap instruksi dalam \texttt{Dockerfile} — seperti \texttt{FROM}, \texttt{RUN}, \texttt{COPY}, dan \texttt{CMD} — menghasilkan satu lapisan (\textit{layer}) \textit{read-only} yang ditumpuk secara berurutan. Docker menyimpan setiap lapisan dalam \textit{cache} lokal dan akan menggunakannya kembali selama tidak ada perubahan pada instruksi tersebut maupun pada lapisan-lapisan sebelumnya [\ref{dockerbp}].

Mekanisme \textit{layer caching} ini memiliki implikasi strategis terhadap urutan penulisan instruksi. Apabila satu lapisan mengalami \textit{invalidation} (misalnya karena perubahan berkas sumber), seluruh lapisan setelahnya akan ikut di-\textit{rebuild} meskipun instruksinya tidak berubah [\ref{dockerbp}]. Oleh karena itu, praktik terbaik (\textit{best practice}) yang direkomendasikan adalah menempatkan instruksi yang jarang berubah (seperti instalasi dependensi sistem) di bagian atas \texttt{Dockerfile}, dan instruksi yang sering berubah (seperti penyalinan kode sumber) di bagian bawah. Secara khusus untuk aplikasi berbasis \textit{package manager} (Composer, npm), teknik menyalin berkas \textit{lock file} terlebih dahulu sebelum menyalin seluruh kode sumber terbukti efektif dalam mempertahankan \textit{cache} instalasi dependensi [\ref{dockerbp}].

\begin{figure}[H]
    \centering
    \includegraphics[width=0.7\textwidth]{layer-caching.png}
    \caption{Ilustrasi mekanisme layer caching pada proses pembangunan Docker image}
\end{figure}

\clearpage
\section{Infrastruktur Pendukung}

\subsection{Containerization Aplikasi Backend Laravel}
Laravel adalah \textit{framework} PHP yang mengadopsi pola arsitektur \textit{Model-View-Controller} (MVC) dan menyediakan fitur bawaan seperti \textit{routing}, \textit{ORM} (Eloquent), \textit{migration} basis data, dan \textit{artisan CLI} [\ref{laravel}]. Dalam konteks \textit{containerization}, aplikasi Laravel dapat dikemas ke dalam \textit{Docker image} dengan menggunakan \textit{base image} PHP yang dilengkapi ekstensi-ekstensi yang diperlukan (\texttt{pdo\_mysql}, \texttt{pdo\_sqlite}). Proses \textit{build} mencakup instalasi dependensi melalui Composer, pembangkitan kunci aplikasi (\texttt{php artisan key:generate}), dan konfigurasi \textit{web server} internal Artisan pada \textit{port} yang di-\textit{expose}.

Fitur \textit{port mapping} pada Docker (\texttt{-p host:container}) memungkinkan layanan di dalam kontainer diakses dari mesin \textit{host} tanpa modifikasi konfigurasi jaringan yang kompleks [\ref{dockerdocs}]. Berkas \texttt{.dockerignore} berperan penting dalam mengecualikan direktori-direktori yang tidak relevan (seperti \texttt{vendor/}, \texttt{node\_modules/}, \texttt{.git/}) dari \textit{build context}, sehingga mempercepat proses transfer dan mengurangi ukuran \textit{image} akhir secara signifikan.

\begin{figure}[H]
    \centering
    \includegraphics[width=0.7\textwidth]{laravel-docker.jpg}
    \caption{Alur containerization aplikasi Laravel menggunakan Docker}
\end{figure}


\chapter{Hasil dan Pembahasan}

\section{Konfigurasi Docker untuk Aplikasi Laravel}

\subsection{Penyusunan Dockerfile}
Langkah fundamental dalam proses \textit{containerization} adalah penyusunan berkas \texttt{Dockerfile} yang mendefinisikan seluruh instruksi pembangunan \textit{image}. Pada praktikum ini, digunakan pendekatan \textit{single-stage build} berbasis \textit{image} dasar \texttt{php:8.3-cli}. Arsitektur \texttt{Dockerfile} dirancang dengan memperhatikan strategi \textit{layer caching} secara eksplisit: file \texttt{composer.json} dan \texttt{composer.lock} disalin terlebih dahulu sebelum seluruh kode sumber aplikasi agar lapisan instalasi dependensi Composer dapat di-\textit{cache} oleh Docker selama tidak ada perubahan pada kedua file tersebut. Instruksi \texttt{RUN} dikonsolidasikan untuk meminimalkan jumlah lapisan, mencakup instalasi dependensi sistem (seperti \texttt{git}, \texttt{unzip}, \texttt{libzip-dev}, \texttt{sqlite3}), ekstensi PHP (\texttt{pdo\_mysql}, \texttt{pdo\_sqlite}, \texttt{zip}), serta pembersihan \textit{cache} \texttt{apt}. Pada tahap akhir, \textit{image} dikonfigurasi untuk mengekspos \textit{port} 8000 dan menjalankan perintah \texttt{php artisan serve} sebagai proses utama kontainer.

\begin{figure}[H]
    \centering
    \includegraphics[width=0.6\textwidth]{1.png}
    \caption{Isi lengkap berkas Dockerfile untuk containerization aplikasi Laravel}
\end{figure}

\subsection{Konfigurasi .dockerignore}
Berkas \texttt{.dockerignore} berfungsi sebagai mekanisme penyaringan (\textit{filtering}) yang mencegah berkas-berkas yang tidak relevan dari proses penyalinan ke dalam \textit{build context} Docker. Konfigurasi yang diterapkan mengecualikan direktori-direktori besar seperti \texttt{/node\_modules}, \texttt{/vendor}, dan \texttt{/.git} yang berpotensi memperbesar ukuran \textit{context} secara signifikan. Selain itu, berkas sensitif seperti \texttt{/.env} juga dikecualikan demi keamanan kredensial. Direktori \textit{cache} dan \textit{session} dari \textit{framework} Laravel (\texttt{/storage/framework/cache/*}, \texttt{/storage/framework/sessions/*}, \texttt{/storage/framework/views/*}) turut diabaikan karena bersifat \textit{runtime-generated} dan tidak diperlukan pada saat proses \textit{build}. Konfigurasi ini secara langsung berdampak pada pengurangan ukuran \textit{transfer context} yang terlihat pada log \textit{build}.

\begin{figure}[H]
    \centering
    \includegraphics[width=0.5\textwidth]{2.png}
    \caption{Konfigurasi berkas .dockerignore untuk penyaringan build context}
\end{figure}

\clearpage
\section{Proses Build Image dan Demonstrasi Layer Caching}

\subsection{Build Pertama (Cold Build)}
Eksekusi perintah \texttt{docker build -t evolusi-api:tugas4 .} untuk pertama kalinya menghasilkan proses \textit{cold build} yang memakan waktu total 88,7 detik (15 langkah). Pada kondisi ini, seluruh instruksi dalam \texttt{Dockerfile} dieksekusi secara penuh tanpa \textit{cache} karena belum ada lapisan (\textit{layer}) yang tersimpan sebelumnya. Proses yang paling memakan waktu adalah pengunduhan dan ekstraksi \textit{base image} \texttt{php:8.3-cli} serta \textit{image} \texttt{composer:latest}, yang secara kumulatif membutuhkan waktu puluhan detik untuk mengambil seluruh \textit{layer} dari \textit{registry} Docker Hub.

\begin{figure}[H]
    \centering
    \includegraphics[width=0.8\textwidth]{3.png}
    \caption{Log eksekusi docker build pertama (cold build) dengan durasi 88,7 detik}
\end{figure}

\subsection{Build Kedua (Full Cache Hit)}
Eksekusi perintah \texttt{docker build} yang kedua tanpa perubahan apapun pada kode sumber menghasilkan waktu \textit{build} yang sangat drastis berkurang menjadi hanya 1,6 detik. Seluruh 8 tahapan (\textit{stage-0}) menampilkan status \texttt{CACHED}, yang membuktikan bahwa Docker berhasil mengenali bahwa tidak ada instruksi yang mengalami perubahan \textit{input} sehingga seluruh lapisan digunakan ulang dari \textit{cache} lokal. Fenomena ini mendemonstrasikan efektivitas arsitektur \textit{layer caching} Docker secara empiris.

\begin{figure}[H]
    \centering
    \includegraphics[width=0.8\textwidth]{4.png}
    \caption{Log eksekusi docker build kedua menunjukkan seluruh layer berstatus CACHED (1,6 detik)}
\end{figure}

\subsection{Build Ketiga (Partial Cache Invalidation)}
Untuk menguji mekanisme \textit{cache invalidation}, dilakukan modifikasi minor pada berkas \texttt{README} (penambahan satu karakter). Eksekusi \textit{build} ketiga memakan waktu 2,7 detik — lebih lama dari \textit{build} kedua namun jauh lebih cepat dari \textit{build} pertama. Analisis log menunjukkan bahwa lapisan 2 hingga 6 (\texttt{apt-get install}, \texttt{COPY composer}, \texttt{WORKDIR}, \texttt{COPY composer.json}, dan \texttt{composer install}) tetap berstatus \texttt{CACHED} karena tidak terpengaruh oleh perubahan tersebut. Namun lapisan 7 (\texttt{COPY . .}) dan 8 (\texttt{RUN cp .env.example}) harus dieksekusi ulang karena perubahan pada berkas sumber meng-\textit{invalidate} konteks direktori. Hasil ini memvalidasi bahwa strategi penyalinan \texttt{composer.json} terlebih dahulu pada \texttt{Dockerfile} berhasil mempertahankan \textit{cache} instalasi dependensi.

\begin{figure}[H]
    \centering
    \includegraphics[width=0.8\textwidth]{5.png}
    \caption{Log eksekusi docker build ketiga pasca edit README, menunjukkan partial cache invalidation (2,7 detik)}
\end{figure}

\subsection{Verifikasi Log Build di Docker Desktop}
Verifikasi melalui antarmuka Docker Desktop pada menu \textit{Builds} mengonfirmasi rekaman ketiga sesi \textit{build} secara kronologis. Terlihat bahwa durasi masing-masing \textit{build} tercatat konsisten dengan observasi pada terminal: \textit{build} pertama berdurasi 1 menit 27 detik, \textit{build} kedua 1,4 detik, dan \textit{build} ketiga 2,4 detik. Seluruh \textit{build} menggunakan \textit{builder} \texttt{desktop-linux} dan berstatus sukses (ditandai ikon centang hijau).

\begin{figure}[H]
    \centering
    \includegraphics[width=0.9\textwidth]{6.png}
    \caption{Tampilan Build history pada Docker Desktop merekam tiga sesi build secara kronologis}
\end{figure}

\clearpage
\section{Eksekusi Container dan Verifikasi Fungsionalitas}

\subsection{Menjalankan Container dan Verifikasi Status}
Setelah \textit{image} berhasil dibangun, kontainer diinisiasi menggunakan perintah \texttt{docker run -d --name laravel-api -p 8000:8000 evolusi-api:tugas4}. Opsi \texttt{-d} menginstruksikan Docker untuk menjalankan kontainer dalam mode \textit{detached} (latar belakang), sedangkan \texttt{-p 8000:8000} memetakan \textit{port} 8000 pada mesin \textit{host} ke \textit{port} 8000 di dalam kontainer. Verifikasi melalui perintah \texttt{docker ps} mengonfirmasi bahwa kontainer \texttt{laravel-api} berstatus \texttt{Up} dengan \textit{port mapping} yang sesuai.

\begin{figure}[H]
    \centering
    \includegraphics[width=0.8\textwidth]{7.png}
    \caption{Eksekusi docker run dan verifikasi status kontainer aktif via docker ps}

    \vspace{0.5cm}

    \includegraphics[width=0.9\textwidth]{8.png}
    \caption{Tampilan kontainer laravel-api berstatus aktif pada antarmuka Docker Desktop}
\end{figure}

\subsection{Verifikasi Halaman Laravel}
Pengujian aksesibilitas dilakukan dengan membuka alamat \texttt{localhost:8000} melalui peramban. Halaman \textit{welcome} Laravel versi 13.33.0 berhasil ter-\textit{render} secara utuh, yang membuktikan bahwa \textit{web server} Artisan di dalam kontainer berfungsi dengan baik dan \textit{port mapping} antara \textit{host} dan kontainer telah terkonfigurasi secara tepat.

\begin{figure}[H]
    \centering
    \includegraphics[width=0.8\textwidth]{9.png}
    \caption{Halaman welcome Laravel v13.33.0 berhasil diakses melalui localhost:8000}
\end{figure}

\subsection{Pengujian API via Postman}
Pengujian lebih lanjut dilakukan pada \textit{endpoint} API \texttt{http://localhost:8000/\allowbreak api/portfolios} menggunakan Postman. Respons yang diterima berstatus \texttt{200 OK} dengan waktu respons 79 ms dan ukuran \textit{payload} 674 B. Data yang dikembalikan berupa larik (\textit{array}) JSON berisi objek portofolio dengan atribut \texttt{id}, \texttt{title}, \texttt{description}, \texttt{image\_url}, dan \texttt{category}. Hasil ini memverifikasi bahwa seluruh \textit{stack} aplikasi Laravel — mulai dari \textit{routing}, \textit{controller}, hingga \textit{database seeder} — beroperasi secara fungsional di dalam lingkungan kontainer Docker.

\begin{figure}[H]
    \centering
    \includegraphics[width=0.9\textwidth]{10.png}
    \caption{Pengujian endpoint API /api/portfolios via Postman mengembalikan respons 200 OK}
\end{figure}

\clearpage
\section{Manajemen Repositori dan Eksekusi Pipeline CI/CD}

Implementasi pada sisi manajemen repositori melibatkan dua lingkungan organisasi GitHub secara simultan. Repositori resmi organisasi mata kuliah (\texttt{KEPL2026}) sedang berstatus \textit{flag} (dalam proses aju banding), yang mengakibatkan fitur GitHub Actions dinonaktifkan secara paksa. Oleh karena itu, eksekusi \textit{pipeline} CI/CD dialihkan ke repositori pada organisasi pribadi (\texttt{AcademyCity-RND}), sementara \textit{push} kode tetap dilakukan ke kedua repositori.

\subsection{Pull Request Feature ke Dev}
Alur kolaborasi dimulai dengan pembuatan \textit{Pull Request} bertajuk \textit{``feat: Integrasi Docker Backend (Tugas 4)''} dari \textit{branch} \texttt{feature/docker-tugas4} menuju \texttt{dev}. Kedua repositori menunjukkan riwayat \textit{commit} yang mencakup penambahan \texttt{Dockerfile} dan \texttt{.dockerignore}, optimasi \textit{layer caching}, penyesuaian versi \textit{base image}, serta perbaikan konfigurasi \texttt{.env}. Status \textit{Pull Request} pada kedua repositori adalah \textit{Ready to merge} tanpa konflik.

\begin{figure}[H]
    \centering
    \includegraphics[width=0.9\textwidth]{11.png}
    \caption{Pull Request feat: Integrasi Docker Backend dari branch feature/docker-tugas4 ke dev (Kiri: repo pribadi AcademyCity-RND, Kanan: repo KEPL2026)}
\end{figure}

\subsection{Pull Request Dev ke Main}
Promosi tahap akhir dilakukan melalui \textit{Pull Request} bertajuk \textit{``Release: Tugas 4 (Docker \& Containerization)''} dari \textit{branch} \texttt{dev} menuju \texttt{main}. Pada repositori KEPL2026, terlihat status \textit{``No conflicts with base branch''} yang mengindikasikan kelancaran proses \textit{merge}. Sementara pada repositori pribadi, notifikasi \textit{``This branch has not been deployed''} terlihat karena proses \textit{deployment} memerlukan persetujuan \textit{reviewer} terlebih dahulu.

\begin{figure}[H]
    \centering
    \includegraphics[width=0.9\textwidth]{12.png}
    \caption{Pull Request Release: Tugas 4 dari branch dev ke main (Kiri: repo pribadi, Kanan: repo KEPL2026)}
\end{figure}

\subsection{Eksekusi Pipeline CI/CD dan Environment Protection}
Pada repositori pribadi yang memiliki GitHub Actions aktif, \textit{pipeline} CI/CD dieksekusi secara otomatis setelah kode masuk ke \textit{branch} \texttt{main}. Tampilan awal menunjukkan 7 \textit{checks} berhasil (\textit{Lint Code}, \textit{Unit Tests}, \textit{Build Application}, \textit{Deploy to Staging}, dan lainnya) dengan 1 \textit{check} berstatus \textit{Waiting} — yaitu \textit{``Deploy to Production''} yang membutuhkan persetujuan manual dari \textit{reviewer} sebagai mekanisme \textit{Environment Protection}. Setelah \textit{reviewer} memberikan persetujuan (\textit{approval}), seluruh 8 \textit{checks} berhasil dilalui dan \textit{deployment} ke lingkungan produksi dieksekusi dengan sukses.

\begin{figure}[H]
    \centering
    \includegraphics[width=0.9\textwidth]{13.png}
    \caption{Pipeline CI/CD pada branch main: 7 checks sukses dengan 1 check (Deploy to Production) menunggu persetujuan reviewer}
\end{figure}

\begin{figure}[H]
    \centering
    \includegraphics[width=0.9\textwidth]{14.png}
    \caption{Seluruh 8 checks berhasil dilalui setelah reviewer memberikan persetujuan deployment ke produksi}
\end{figure}

\section{Tautan Repositori Kode Sumber}

Seluruh kode sumber praktikum ini tersedia pada dua repositori berikut:
\begin{itemize}
    \item Repo Matkul KEPL2026:\\
    \url{https://github.com/KEPL2026/evolusi-pl-24-542159-SV-25009-v1}

    \item Repo Pribadi (AcademyCity-RND):\\
    \url{https://github.com/AcademyCity-RND/evolusi-pl-24-542159-SV-25009-v2}
\end{itemize}


\chapter{Kesimpulan}
Berdasarkan praktikum mengenai pengenalan Docker dan \textit{containerization} pada aplikasi \textit{backend} Laravel yang telah dilakukan, dapat ditarik beberapa kesimpulan utama sebagai berikut:
\begin{enumerate}
    \item Penyusunan \texttt{Dockerfile} dengan pendekatan \textit{single-stage build} berbasis \texttt{php:8.3-cli} berhasil mengemas seluruh lingkungan aplikasi Laravel — termasuk dependensi sistem, ekstensi PHP, dan \textit{package} Composer — ke dalam satu \textit{Docker image} yang portabel dan dapat direproduksi secara konsisten di berbagai mesin.
    \item Strategi \textit{layer caching} yang diterapkan melalui pemisahan penyalinan berkas \texttt{composer.json} dan \texttt{composer.lock} sebelum kode sumber terbukti sangat efektif secara empiris: \textit{build} pertama memakan waktu 88,7 detik (\textit{cold build}), \textit{build} kedua hanya 1,6 detik (\textit{full cache hit}), dan \textit{build} ketiga 2,7 detik (\textit{partial cache invalidation} pasca modifikasi minor). Hal ini mendemonstrasikan pentingnya urutan instruksi dalam \texttt{Dockerfile} terhadap efisiensi proses pembangunan.
    \item Konfigurasi \texttt{.dockerignore} yang mengecualikan direktori \texttt{vendor/}, \texttt{node\_modules/}, \texttt{.git/}, dan berkas \texttt{.env} berperan penting dalam memperkecil ukuran \textit{build context} dan mencegah tereksposnya kredensial sensitif ke dalam \textit{image}.
    \item Mekanisme \textit{port mapping} (\texttt{-p 8000:8000}) berhasil menghubungkan layanan Artisan di dalam kontainer dengan mesin \textit{host}, yang diverifikasi melalui aksesibilitas halaman Laravel pada peramban serta keberhasilan pengujian \textit{endpoint} API via Postman dengan respons \texttt{200 OK}.
    \item Alur manajemen kode melalui \textit{Pull Request} bertahap (\texttt{feature} → \texttt{dev} → \texttt{main}) yang diintegrasikan dengan \textit{pipeline} CI/CD dan \textit{Environment Protection} menunjukkan penerapan praktik terbaik dalam \textit{DevOps}, di mana kode hanya dapat mencapai lingkungan produksi setelah melewati seluruh \textit{quality gates} dan memperoleh persetujuan manual dari \textit{reviewer}.
\end{enumerate}


\chapter{Daftar Pustaka}

\begin{enumerate}
    \item \label{felter} W. Felter, A. Ferreira, R. Rajamony, \& J. Rubio. \textit{An Updated Performance Comparison of Virtual Machines and Linux Containers}. IBM Research Division, 2015. [Online]. Tersedia: \url{https://dominoweb.draco.res.ibm.com/reports/rc25482.pdf}

    \item \label{dockerdocs} Docker, Inc. \textit{Docker Documentation: Docker Overview}. Docker Docs, 2024. [Online]. Tersedia: \url{https://docs.docker.com/get-started/docker-overview/}

    \item \label{dockerbp} Docker, Inc. \textit{Dockerfile Best Practices: Leverage Build Cache}. Docker Docs, 2024. [Online]. Tersedia: \url{https://docs.docker.com/build/cache/}

    \item \label{laravel} Laravel LLC. \textit{Laravel Documentation}. Laravel, 2024. [Online]. Tersedia: \url{https://laravel.com/docs}

    \item \label{sourcecode} Source Code Praktikum:
    \begin{itemize}
        \item Repo Matkul KEPL2026:\\
        \url{https://github.com/KEPL2026/evolusi-pl-24-542159-SV-25009-v1}

        \item Repo Pribadi:\\
        \url{https://github.com/AcademyCity-RND/evolusi-pl-24-542159-SV-25009-v2}
    \end{itemize}
\end{enumerate}